\documentclass{article}
\usepackage{amsmath, amssymb}
\usepackage{cancel}
\usepackage{tikz}
\usepackage{mdframed}
\usepackage{hyperref}
\usetikzlibrary{arrows.meta, decorations.markings}

\begin{document}

\section{Lecture 14: Plan for Class}
Last time:
\begin{enumerate}
    \item Averaging \S 4.1-2
    \item Duffing equation
\end{enumerate}

Today:
\begin{enumerate}
    \item Project information
\end{enumerate}

\section{Averaging}
Suppose that we have a system 
\[
\dot{x} = \epsilon f(x,t,\epsilon)
\]
where $x\in\mathbb{R}^n$ and $f$ is $T$-periodic. The \textit{averaged system} is denoted 
\[
\dot{y}=\epsilon \bar{f}(y)
\]
where $\bar{f}$ is 
\[
\bar{f}(x)=\int^T_0 f(x,t,0)dt
\]
and $x$ is fixed. Note that the phase portrait of $\dot{y}=\epsilon \bar{f}(y)$ looks a lot like the phase portrait of the Poincare map of $\dot{x} = \epsilon f(x,t,\epsilon)$.

\subsection{Duffing equation}

The duffing equation can be written as 
\[
\ddot{x} + \omega_0^2 x + \epsilon \delta \dot{x} + \epsilon \alpha x^3 = \epsilon \gamma \cos \omega t
\]
where $\epsilon$ is small and $\omega$ is close to $\omega_0$. We can rescale $t$ to get 
\[
\ddot{x} + x + \delta \dot{x} + \alpha x^3 = \gamma \cos \omega t
\]
where $\delta,\alpha,\gamma$ are small and $\omega$ is close to 1. We can further analyze this system via using rotating coordinates via 
\[
\begin{bmatrix}
    u \\ v
\end{bmatrix} = A \begin{bmatrix}
    x \\ \dot{x}
\end{bmatrix}
\]
where $A$ is 
\[
A = \begin{bmatrix}
    \cos \omega t & -\frac{1}{\omega}\sin \omega t \\
    -\sin \omega t & -\frac{1}{\omega} \cos \omega t
\end{bmatrix}
\]
We can further reduce this with polar coordinates
\begin{align*}
    u &= r\cos \varphi \\
    v &= r\sin\varphi
\end{align*}
to get 
\begin{align*}
    \dot{r} &= \frac{1}{2\omega}\left[-\omega \delta r - \gamma \sin\varphi\right]\\
    r\dot{\varphi} &= \frac{1}{2\omega}\left[\left(1-\omega^2\right)r + \frac{3\alpha}{4}r^3-\gamma \cos\varphi \right]
\end{align*}

\section{Summary of the other averaging results}
\begin{enumerate}
    \item[\S 4.3] bifurcations
    \begin{itemize}
        \item local bifurcations of averaged system also tell you about local bifurcations of original system
    \end{itemize}
    \item[\S 4.4] global behavior - WATCH OUT
    \begin{itemize}
        \item hyperbolic features preserved: hyperbolic periodic orbits of the averaged equations imply hyperbolic invariant torus of original equations
        \item non-hyperbolic features not preserved: example is the Duffing equation with $\delta=0$ where the system is Hamiltonian (this comes from some fragile features in the system)
    \end{itemize}
\end{enumerate}

\subsection{Homoclinic tangles}

We can get homoclinic tangles in Poincare maps, but not in flows! Therefore, the averaging method can't capture homoclinic tangles.

\section{Melnikov's method \S 4.5}

Melnikov's method is used for time-periodic perturbations of planar ($\mathbb{R}^2$) Hamiltonian systems, but can be extended for many other systems. But let's say we have s system 
\[
\dot{x} = f(x) + \epsilon g(x,t)
\]
where $f$ is Hamiltonian
\begin{align*}
    f_1 &= \frac{\partial H}{\partial x_2}\\
    f_2 &= -\frac{\partial H}{\partial x_1}
\end{align*}
and $g$ is of the form 
\[
g(x,t+T) = g(x,t)
\]
for all $x,t$. 

\subsection{Duffing equation}

Our motivating example is the duffing equation with a minus sign 
\[
\ddot{x} - x + x^3 = \epsilon\left(-\delta \dot{x} + \gamma \cos \omega t \right)
\]
An unperturbed system is Hamiltonian, and we could use the energy function 
\[
V(x) = \frac{x^4}{4}-\frac{x^2}{2}
\]
where the Hamiltonian becomes 
\[
H(x,y) = \frac{y^2}{2} + V(x)
\]
where 
\[
\dot{x} = \frac{\partial H}{\partial y}=y
\]
and 
\[
\ddot{x} = \dot{y} = -\frac{\partial H}{\partial x}=-V'(x) = -x^3 + x
\]

\subsubsection{Phase portraits}

Rowley is now drawing some phase portraits that describe how the Hamiltonian changes when $\epsilon\neq 0$. 

What if $\delta>0$ and $\gamma >0$? Consider the Poincare map. Can stable and unstable manifolds intersect?

\subsection{Melnikov's method}

Construct the Melnikov function 
\[
M\left(t_0\right) = \int_{-\infty}^{\infty} \left(f\times g\right)\left(q(t),t+t_0\right)dt
\]
where $q(t)$ is the homoclinic orbit of the unperturbed system. If $M(t_0)$ has simple zeros $t_0\in \left[0,T\right]$, then there exists transversal intersections of stable and unstable manifolds. If $M(t_0)$ does not have simple zeros, then there are no transversal intersections.

\subsection{Melnikov function for duffing equation}

For the duffing equation, we have 
\[
f\begin{pmatrix}
    x \\ y
\end{pmatrix}=\begin{pmatrix}
    y \\ x-x^3
\end{pmatrix} \qquad \text{and} \qquad g\left(\begin{pmatrix}
    x \\ y
\end{pmatrix},t\right)=\begin{pmatrix}
    0 \\ -\delta y + \gamma \cos \omega t
\end{pmatrix}
\]
and our $q(t)$ for the homoclinic orbit is 
\[
q(t) = \begin{bmatrix}
    \sqrt{2} sech t \\ -\sqrt{2} sech t \cdot \tanh t
\end{bmatrix}= \begin{bmatrix}
    x(t) \\ y(t)
\end{bmatrix}
\]
Lots of calculus later, and we have 
\[
M(t_0) = \sqrt{2} \gamma \pi \omega sech\left(\frac{\pi \omega}{2}\right)\sin \omega t_0 - \frac{4 \delta }{3}
\]
and we get zeros (and a homoclinic tangle) if 
\[
\frac{\gamma}{\delta} > \frac{4}{3} \frac{\pi \cosh \left(\frac{\pi\omega}{2}\right)}{\sqrt{2} \omega}
\]

\subsection{Numerical example}

Now Rowley is showing us a numerical example.



\end{document}